\newpage
\begin{center}
  \textbf{\large 1. МЕТОДЫ ПЕРЕХВАТА СЕТЕВОГО ТРАФИКА МОБИЛЬНЫХ ПРИЛОЖЕНИЙ ЧЕРЕЗ ПРОКСИ-СЕРВЕР}
\end{center}
\refstepcounter{chapter}
\addcontentsline{toc}{chapter}{1. МЕТОДЫ ПЕРЕХВАТА СЕТЕВОГО ТРАФИКА МОБИЛЬНЫХ ПРИЛОЖЕНИЙ ЧЕРЕЗ ПРОКСИ-СЕРВЕР}


\section{Применимость обратной разработки к анализу мобильных приложений на наличие недокументированных возможностей}
% было порядка 60 строк

Понимание работы современных мобильных приложений подразумевает знание о работе клиентской и серверной части системы. Клиентская часть приложения отвечает за представление, хранение и локальную обработку данных, необходимых для работы приложения, в том время как серверная часть обрабатывает ключевые бизнес функции, контролирует выдачу доступа к личной и персональной информации. Более того, серверные системы контролируют доступ к распространению и скачиванию клиентского приложения, поэтому представляют наибольший интерес с точки зрения анализа на уязвимости и недокументированные возможности.

Традиционный метод, использующийся для анализа клиентских приложений - обратная разработка. Обратная разработка приложения на основе исполняемого файла и конфигураций в виде ключей доступа до сервера с закрытым исходным кодом, временных сертификатов и локального хранилища персональных, публичных и секретных данных позволяет целиком воспроизвести поведение системы на локальном устройстве с архитектурой, поддерживаемой целевым клиентским приложением. Однако, полноценная обратная разработка приложения - долгий и тяжелый процесс, который зависит от платформы, для которой приложение предназначено, наличия в системе динамических библиотек, реализующих необходимые функции для обеспечения работоспособности клиентского кода. Успех обратной разработки может зависеть от доступных вычислительных мощностей для раскрытия зашифрованной информации, распространяющейся совместно с открытой частью приложения.

Обратная разработка, будучи известным методом для анализа приложений на защищенность и недокументированные возможности, учитывается как доступный вектор атаки на клиентское приложение с самого начала разработки мобильного приложения, и потому разработчики и поставщики клиентских систем заранее предусматривают защитные функции, усложняющие процесс обратной разработки мобильных приложений. Самая распространенная такая функция - удаление из распространяемого файла отладочной информации, обфускация исполняемого кода и использование нестандартных операций при работе с памятью и обработкой данных. Данные защитные меры просты во внедрении в процессы сборки и передачи для уже готовых мобильных приложений, и при этом достаточно эффективны в повышении запутанности кода исполняемого файла, что существенно усложняет процесс обратной разработки.

Осведомлённость о рисках обратной разработки и наличие публично доступных средств от её защиты с открытым исходным кодом усложняют и замедляют процесс анализа приложения. Тем временем, для мобильных приложений, распространяемых по сети, время, затраченное на анализ, является ключевой метрикой указывающей на эффективность выбранного метода для анализа мобильного приложения на защищённость, так как новая версия системы сразу же после загрузки в централизованные репозитории мобильных приложений становится доступной для установки на все поддерживаемые устройства. В современных командах, придерживающихся гибкой итеративной разработки, релизный цикл может составлять сутки и даже меньшее число времени. В таких условиях необходимо так же быстро проверять систему на наличие уязвимостей и недокументированных возможностей. И если эта проверка занимает больше времени, чем будет затрачено на выпуск новой версии исследуемого приложения с закрытым исходным кодом, то проверяющей стороне потребуется одновременно вести анализ для нескольких версий одного и того же приложения.

Другой немаловажный фактор, влияющий на выбор обратной разработки в качестве предпочтительного метода для исследования приложения на соответствие стандартам безопасности, заключается в сложности переносимости метода между различными приложениями. Так, для двух версий одного и того же мобильного приложения, поддерживающих разные архитектуры компьютера, требуется изменять метод обратной разработки и проводить двойную работу по интерпретации кода исполняемого файла и пользовательских файлов, создаваемых в процессе работы приложения или поставляющихся вместе с основной программой.

При всех своих недостатках, обратная разработка всё же имеет и преимущества. По результатам метода можно полностью понимать работу клиентской части приложения и делать обоснованные предположения о серверной составляющей. Это позволяет обозначить набор тестов, гарантирующих с высокой долей вероятности, стабильную и безопасную обработку, хранение и передачу пользовательских данных внутри контура системы и за её пределами. Формальные шаги для верификации корректной работы мобильного приложения можно использовать для верификации исследования и составления отчётов о безопасности использования того или иного приложения в ситуациях с критической ценностью персональных или секретных данных.

Обратная разработка, проведённая в полном объёме, позволяет воссоздать исследуемое приложение и реализовать его аналог с открытым исходным кодом. Сравнение исходного приложения и его двойника с открытым исходным кодом может быть выполнено быстрее и проще чем полный анализ приложения с закрытым исходным кодом, и в том числе может быть выполнено автоматически в ряде случаев. Поэтому возможность реализации самостоятельной копии является ключевым конкуретным преимуществом данного метода.

На основе имеющихся преимуществ и недостатков были сформулированы цели на поиск нового метода, имеющего достаточную степень точности в определении недокументированных возможностей и уязвимостей мобильных приложений, и одновременно доступным для апробации в условиях ограниченных ресурсов времени, вычислительных мощностей и доступа к устройствам с архитектурой, поддерживаемой целевым приложением.

\section{Проведение автоматического сканирования на наличие уязвимостей мобильных приложений}
% было порядка 90 строк



\section{Цели и задачи работы}

\textbf{Цель работы} -- установить связь дальнодействия притяжения потенциала взаимодействия и спектров возбуждений с транспортными свойствами жидкостей, а также влияние на скорость нуклеации.

\textbf{Задачи работы:}
\begin{enumerate}
\item Расчет фазовых диаграмм для 2D и 3D систем частиц, взаимодействующих посредством обобщенного потенциала Леннарда-Джонса с различными степенями притяжения.
\item Адаптация метода кластеризации данных DBSCAN для изучения молекулярных систем и его сравнение с другими методами.
\item Расчет и анализ транспортных свойств и коллективных возбуждений на жидкостных бинодалях.
\item Применение нового метода распознавания фаз для изучения скорости нуклеации в переохлажденных системах Леннарда-Джонса с различным дальнодействием притяжения.
\end{enumerate}
